\section{Technical Background}
\label{app:technical-background}

\subsection{Quantum states, replacement channels, and trace norm}

An \(n\)-qubit density operator is a positive semidefinite matrix of trace one
on \((\mathbb C^2)^{\otimes n}\). The trace norm is
\(\traceNorm{A}=\Tr\sqrt{A^\dagger A}\), and the operator norm is
\(\opNorm{A}\). We repeatedly use
\[
    \abs{\Tr[AB]}
    \leq\opNorm{A}\traceNorm{B},
\]
unitary invariance of the trace norm, and the fact that two density operators
have trace-norm distance at most two.

For completeness, the replacement map \(\mathcal R_\sigma(A)=\Tr[A]\sigma\)
is completely positive and trace preserving; this is a finite-dimensional
instance of the operator-sum representation of quantum operations
\citep{Kraus1971}. Indeed, if
\(\sigma=\sum_a r_a\lvert a\rangle\!\langle a\rvert\) and
\(\{\lvert b\rangle\}_b\) is an orthonormal basis, the Kraus operators
\(K_{ab}=\sqrt{r_a}\lvert a\rangle\!\langle b\rvert\) satisfy
\[
    \sum_{a,b}K_{ab}A K_{ab}^\dagger=\Tr[A]\sigma,
    \qquad
    \sum_{a,b}K_{ab}^\dagger K_{ab}=I.
\]
Thus the reset update in \cref{eq:reset-channel} is a convex combination of
two quantum channels.

\subsection{Absolute regularity and the independent-block theorem}

For sigma-algebras \(\mathcal A\) and \(\mathcal B\), their absolute-regularity
coefficient can be written as
\[
    \beta(\mathcal A,\mathcal B)
    =
    \frac12
    \sup_{\{A_i\},\{B_j\}}
    \sum_{i,j}
    \abs{\Prob(A_i\cap B_j)-\Prob(A_i)\Prob(B_j)},
\]
where the supremum ranges over finite measurable partitions from the two
sigma-algebras. The coefficient is monotone under inclusion of either
sigma-algebra, and measurable transformations cannot increase it
\citep{Bradley2005}.

We use the following direct specialization of
\citet[Theorem~2]{MohriRostamizadeh2008}. Let \((Q_j)_{j\in\Z}\) be stationary
and beta mixing, let \(\mathcal F\) be a fixed measurable class of functions
into \([0,B]\), and let \(m=2\mu a\). Divide \(Q_1,\ldots,Q_m\) into
\(2\mu\) consecutive blocks of length \(a\), and form \(\cS_\mu\) by
selecting the same fixed coordinate from every odd block. With the source
paper's empirical Rademacher convention
\begin{equation}
    \hRad_{\cS_\mu}(\mathcal F)
    :=
    \E_{\bm\sigma}\!\left[
      \sup_{f\in\mathcal F}
      \abs{\frac2\mu\sum_{j=1}^{\mu}
      \sigma_j f(\widetilde Q_j)}
    \right],
    \label{eq:source-rademacher-definition}
\end{equation}
if
\[
    \delta_Q'
    :=
    \delta-4(\mu-1)\beta_Q(a)>0,
\]
then, with probability at least \(1-\delta\), simultaneously for every
\(f\in\mathcal F\),
\begin{equation}
    \E[f(Q_0)]
    \leq
    \frac1m\sum_{j=1}^m f(Q_j)
    +\hRad_{\cS_\mu}(\mathcal F)
    +3B\sqrt{\frac{\log(4/\delta_Q')}{2\mu}}.
    \label{eq:mr-background}
\end{equation}
The theorem is one-sided. A two-sided statement would require an additional
reverse-deviation argument and a corresponding confidence allocation. We use
\cref{eq:source-rademacher-definition} throughout so that the imported theorem
and all constants in our specialization share exactly the same convention.

\subsection{Spectral representation of the Mat\'ern kernel}
\label{app:matern-spectral}

By Bochner's theorem \citep{Bochner1933}, the normalized stationary
Mat\'ern kernel has a probability spectral measure:
\begin{equation}
    \kappa(z,z')
    =
    \E_\Omega\!\left[e^{\ii\Omega^\top(z-z')}\right]
    =
    \E_\Omega\!\left[\cos(\Omega^\top(z-z'))\right].
    \label{eq:matern-bochner}
\end{equation}
For the parameterization in \cref{eq:matern-profile}, its density on
\(\R^p\) is
\begin{equation}
    p_{\nu,\xi}(\omega)
    =
    \frac{\Gamma(\nu+p/2)\,\xi^p}
         {\Gamma(\nu)(2\nu\pi)^{p/2}}
    \left(1+\frac{\xi^2\norm{\omega}_2^2}{2\nu}\right)^{-(\nu+p/2)},
    \label{eq:matern-spectral-density}
\end{equation}
which is a centered multivariate Student density with \(2\nu\) degrees of
freedom and scale matrix \(\xi^{-2}I\); see
\citet{Genton2001,PorcuEtAl2024}. It is strictly positive everywhere. For
\(\nu>1\), its covariance exists and is
\begin{equation}
    \E[\Omega\Omega^\top]
    =
    \frac{\nu}{(\nu-1)\xi^2}I.
    \label{eq:matern-spectral-covariance}
\end{equation}
These identities yield the section-stability constant used in
\cref{lem:matern-stability}. Positivity of the density also yields strict
positive definiteness on distinct Euclidean points, as made explicit in the
proof of \cref{thm:interpolation}; the general RKHS interpolation facts used
there are standard \citep{Aronszajn1950,KimeldorfWahba1971}.

\subsection{Local Pauli classical shadows}

In one snapshot of the local-Pauli shadow protocol, each qubit is measured in
an independently and uniformly chosen \(X\), \(Y\), or \(Z\) basis. Inversion
of the one-qubit measurement channel multiplies a nonidentity Pauli factor by
three. Hence a Pauli string of weight \(q\) has squared shadow norm \(3^q\).
The median-of-means guarantee of \citet{HuangKuengPreskill2020} therefore gives
simultaneous estimation of \(M\) nonidentity, weight-at-most-\(k\) Pauli
expectations with
\(O(3^k\varepsilon^{-2}\log(M/\delta))\) independent snapshots. Identity
observables, if included in a feature ledger, are known exactly and require no
measurement.

\subsection{Provenance of imported ingredients}

The replacement-channel representation follows the standard operator-sum
framework \citep{Kraus1971}, and the spectral representation of the stationary
Mat\'ern kernel invokes Bochner's theorem \citep{Bochner1933}. The finite-set
projection result is the standard JL argument
\citep{DasguptaGupta2003,Woodruff2014}; the dependent-data inequality is
imported from \citet{MohriRostamizadeh2008}; the RKHS and representer facts are
standard consequences of \citet{Aronszajn1950,KimeldorfWahba1971}; and the
finite-union step uses the bounded-difference method of \citet{McDiarmid1989}
and standard Rademacher-complexity tools \citep{BartlettMendelson2002}. The
results specific to \textsc{QuaRK} are the way these ingredients are composed
with its reset-contractive dynamics, finite-window approximation, observable
map, and fixed-design temporal learning problem.
